\documentclass[10pt,a4paper,sans]{moderncv}
\moderncvstyle{classic}
\moderncvcolor{blue}
\usepackage[utf8]{inputenc}
\usepackage[scale=0.8]{geometry}
\usepackage{helvet}

\name{Alexandro}{Disla}
\title{Spécialiste MEAL | Suivi, Évaluation, Redevabilité et Apprentissage}
\address{Rue Saint-Louis Jeanty \#14, Route de Frère}{Port-au-Prince, Haïti}{}
\phone[mobile]{(+509) 4148-3700}
\email{alexandrodisla@hotmail.com}
\extrainfo{GitHub: \url{https://github.com/AD0791}}

\begin{document}
\makecvtitle

\section{Profil Professionnel}
\cvitem{}{Professionnel MEAL avec plus de cinq ans d'expérience sur trois secteurs en Haïti --- santé, eau et assainissement, éducation --- consacrés à la conception des systèmes qui rendent les données d'un programme crédibles : cadres logiques et tableaux de suivi des indicateurs, protocoles d'assurance qualité des données (DQA) établis et supervisés, collecte numérique standardisée afin que la donnée soit juste à la saisie plutôt que réparée ensuite. Économiste appliqué de formation et \textbf{ingénieur de données} en parallèle, je sais construire, sécuriser et administrer les bases et les tableaux de bord dont dépend un bureau pays, et pas seulement les spécifier. Résidant à Port-au-Prince ; français, anglais et créole haïtien courants.}

\section{Compétences Clés}
\cvitem{}{
\begin{itemize}
    \item \textbf{Cadres MEAL \& gestion axée sur les résultats :} Cadres logiques, chaînes de résultats, tableaux de suivi des indicateurs (ITT), plans de suivi, valeurs de référence, cibles et désagrégations exigées.
    \item \textbf{Qualité et protection des données :} Protocoles DQA établis et supervisés, vérification par remontée aux sources primaires et suivi des actions correctives jusqu'à clôture ; administration de bases relationnelles, contrôle d'accès par rôle, stockage sécurisé, dédoublonnage des listes de bénéficiaires.
    \item \textbf{Encadrement, renforcement de capacités \& coordination :} Supervision d'enquêteurs et d'agents de collecte, encadrement technique de personnel de suivi-évaluation, formation aux outils de collecte et à la qualité des données ; travail avec les partenaires gouvernementaux (DINEPA, MSPP) et traduction des constats techniques pour une direction non technique.
\end{itemize}}

\section{Expériences MEAL et Qualité de Programme}

\cventry{Jan 2026 -- Mars 2026}{Data Analyst \& Consultant MEAL}{Anseye Pou Ayiti}{Haïti (Télétravail)}{}{
\begin{itemize}
    \item Définition de la stratégie de collecte et mise en œuvre d'indicateurs de performance pour des programmes de leadership multisectoriels.
    \item Codage des formulaires XLSForm pour ODK et Google Forms avec logiques de saut et contraintes de validation, et maintenance des scripts d'intégration Python sur AWS.
    \item Supervision de la qualité des données remontées du terrain, administration du système ODK/Admin Panel et formation des agents de terrain aux outils de collecte.
\end{itemize}}

\cventry{Oct 2023 -- Jan 2026}{Consultant mWater \& Analyste de Données}{HANWASH}{Télétravail (Upwork)}{}{
\begin{itemize}
    \item \textbf{Conception MEAL :} Raffinement des plans MEAL et des tableaux de suivi des indicateurs (ITT) alignés sur les normes JMP pour des programmes d'eau et d'assainissement.
    \item \textbf{Gestion des données :} Déploiement d'enquêtes géospatiales complexes sur mWater, administration de la base et résolution des problèmes de collecte terrain pour préserver l'intégrité des données.
    \item \textbf{Automatisation \& renforcement de capacités :} Création de tableaux de bord automatisés (mWater, Power BI) connectés aux API de collecte pour un reporting en temps réel ; production de guides techniques et formation des équipes à leur usage.
\end{itemize}}

\cventry{Oct 2021 -- Août 2024}{Officier de Suivi \& Évaluation}{Impact Youth Project, Caris Foundation International}{Haïti}{}{
\begin{itemize}
    \item \textbf{Systèmes de collecte :} Mise en place de systèmes de collecte numérique (CommCare, HIVHaiti) sur des programmes de santé multi-sites, avec standardisation des formulaires et des flux entre sites.
    \item \textbf{Assurance qualité (DQA) :} Établissement et supervision des protocoles de Data Quality Assessment sur les données intégrées MySQL, avec remontée des chiffres rapportés aux registres sources et suivi des actions correctives jusqu'à clôture.
    \item \textbf{Encadrement :} Supervision des agents de collecte sur l'ensemble des sites et encadrement technique du personnel de suivi-évaluation sur les procédures de collecte et les contrôles qualité.
    \item \textbf{Reporting \& tableaux de bord :} Développement de rapports d'activité via Python et SQL, réduisant le temps de traitement manuel de 40 \% ; maintenance des tableaux de bord programme sur AWS (Quarto) et Power BI pour le suivi hebdomadaire des indicateurs.
\end{itemize}}

\cventry{Nov 2015 -- Août 2023}{Économiste -- Direction de Planification Économique et Sociale (DPES)}{Ministère de la Planification et de la Coopération Externe (MPCE)}{Port-au-Prince}{}{Contribution aux rapports de suivi de l'exécution du Plan Triennal d'Investissement 2014--2016 ; affectation à l'unité statistique et modélisation, avec participation à la construction du << Modèle de Planification du Développement >>, outil de prévision, de simulation et d'analyse d'impact des politiques publiques ; suivi des indicateurs macroéconomiques avec le Ministère de l'Économie et des Finances.}

\section{Expériences en Ingénierie de Données et Systèmes}

\cventry{Févr 2026 -- Présent}{Ingénieur Logiciel (Fullstack)}{Tekkod LLC}{Télétravail}{}{Direction d'un programme de modernisation d'infrastructure mobile, avec fiabilité hors ligne et sécurité des données pour un usage terrain ; construction de tableaux de bord Looker Studio avec recherche et filtrage pour la visibilité opérationnelle.}

\cventry{Mars 2021 -- Mai 2024}{Développeur Backend \& Ingénieur de Données}{Tekkod LLC}{Télétravail}{}{Mise en place de pipelines ETL (Apache Beam) migrant les données MongoDB vers BigQuery et administration de bases relationnelles en accès concurrent ; développement de services REST sécurisés (FastAPI, JWT) avec contrôle d'accès par rôle et documentation OpenAPI complète.}

\cventry{Nov 2015 -- Mars 2021}{Consultant Logiciel \& Spécialiste Intégration de Données}{CassionSoft (Projet UNOPS)}{Haïti}{}{Travaux d'intégration de données sur un projet des Nations Unies, avec une exposition directe aux procédures et aux standards documentaires onusiens ; implémentation de systèmes d'authentification RBAC (OAuth2/JWT) et validation de la logique d'intégration par des tests unitaires rigoureux.}

\section{Formation}
\cvitem{2010--2014}{Diplôme d'Études Supérieures en Économie Appliquée --- C.T.P.E.A (Centre de Techniques de Planification et d'Économie Appliquée), Port-au-Prince}
\cvitem{2009--2010}{Baccalauréat (Mention) --- Institution Saint-Louis de Gonzague}
\cvitem{Autoformation}{FlatIron School (Software Engineering), CodeWithMosh, Udemy}

\section{Compétences Techniques}
\cvitem{Outils MEAL}{\textbf{KoboToolbox, ODK, XLSForm}, CommCare, mWater, HIVHaiti, cadres logiques, ITT, protocoles DQA}
\cvitem{Analyse}{\textbf{Excel avancé}, \textbf{Power BI (expert)}, SQL, Python (Pandas), R (RMarkdown/Quarto), Looker Studio, modélisation statistique et échantillonnage}
\cvitem{Infrastructure}{MySQL, PostgreSQL, MongoDB, Azure, AWS, Google Cloud (BigQuery), Docker, Git}
\cvitem{Langues}{\textbf{Français (maternel)}, \textbf{Anglais (courant)}, \textbf{Créole haïtien (maternel)}}

\end{document}
